\documentclass[UTF8,11pt,fontset=none]{ctexart}
\usepackage[a4paper,margin=23mm,headheight=15pt,footskip=11mm]{geometry}
\usepackage{fontspec}
\IfFontExistsTF{TeX Gyre Termes}{\setmainfont{TeX Gyre Termes}}{\setmainfont{lmroman10-regular.otf}[BoldFont=lmroman10-bold.otf,ItalicFont=lmroman10-italic.otf,BoldItalicFont=lmroman10-bolditalic.otf]}
\IfFontExistsTF{TeX Gyre Heros}{\setsansfont{TeX Gyre Heros}}{\setsansfont{lmsans10-regular.otf}[BoldFont=lmsans10-bold.otf,ItalicFont=lmsans10-oblique.otf]}
\IfFontExistsTF{DejaVu Sans Mono}{\setmonofont{DejaVu Sans Mono}[Scale=MatchLowercase]}{\setmonofont{lmmono10-regular.otf}}
\IfFontExistsTF{Noto Serif CJK SC}{\setCJKmainfont{Noto Serif CJK SC}[AutoFakeSlant=0.15]}{\setCJKmainfont{FandolSong-Regular.otf}[BoldFont=FandolSong-Bold.otf,ItalicFont=FandolKai-Regular.otf]}
\IfFontExistsTF{Noto Sans CJK SC}{\setCJKsansfont{Noto Sans CJK SC}}{\setCJKsansfont{FandolHei-Regular.otf}}
\IfFontExistsTF{Noto Sans Mono CJK SC}{\setCJKmonofont{Noto Sans Mono CJK SC}}{\setCJKmonofont{FandolFang-Regular.otf}}
\usepackage{amsmath,amssymb,amsthm,mathtools,mathrsfs,esint,bm}
\usepackage{booktabs,longtable,array,calc}
\usepackage{enumitem}
\setlist{nosep,leftmargin=1.7em,topsep=3pt}
\usepackage{fvextra}
\DefineVerbatimEnvironment{Highlighting}{Verbatim}{breaklines,breakanywhere,fontsize=\small}
\usepackage{xurl}
\usepackage[unicode,hidelinks,bookmarksnumbered]{hyperref}
\usepackage{bookmark}
\usepackage{fancyhdr}
\pagestyle{fancy}\fancyhf{}
\fancyhead[L]{\small 第一科学主线：近期研究全稿}
\fancyhead[R]{\small 2026-10-07}\fancyfoot[C]{\thepage}
\renewcommand{\headrulewidth}{0.3pt}
\providecommand{\tightlist}{\setlength{\itemsep}{1pt}\setlength{\parskip}{0pt}}
\providecommand{\passthrough}[1]{#1}
\newcommand{\Ric}{\operatorname{Ric}}
\newcommand{\Vol}{\operatorname{Vol}}
\newcommand{\Rm}{\operatorname{Rm}}
\newcommand{\supp}{\operatorname{supp}}
\newcommand{\Var}{\operatorname{Var}}
\newcommand{\tr}{\operatorname{tr}}
\newcommand{\avint}{\fint}
\newcommand{\Jcal}{\mathcal J}
\newcommand{\Ecal}{\mathcal E}
\newcommand{\dd}{\,\mathrm d}
\newcommand{\Rea}{\operatorname{Re}}
\setlength{\parindent}{2em}
\setlength{\parskip}{3pt plus 1pt}
\setlength{\emergencystretch}{3em}
\linespread{1.10}\allowdisplaybreaks[2]
\setlength{\abovedisplayskip}{8pt plus 2pt minus 2pt}
\setlength{\belowdisplayskip}{8pt plus 2pt minus 2pt}
\renewcommand{\arraystretch}{1.12}
\setcounter{tocdepth}{1}\setcounter{secnumdepth}{2}
\pdfstringdefDisableCommands{\def\fint{int}\def\mathcal#1{#1}\def\mathfrak#1{#1}\def\mathscr#1{#1}\def\operatorname#1{#1}\def\varepsilon{epsilon}\def\chi{chi}\def\beta{beta}\def\rho{rho}\def\tau{tau}\def\Phi{Phi}\def\sigma{sigma}\def\kappa{kappa}\def\nabla{D}\def\mathbb#1{#1}\def\({}\def\){}\def\textsuperscript#1{#1}\def\textsubscript#1{#1}}
\title{任意复双曲领口中的解析低成本构造}
\author{第一科学主线·近期研究分卷}
\date{2026年10月7日}
\begin{document}
\maketitle
\noindent\textbf{原稿编号：}R09。

用户系统先得到严格次四单元及颈部常模的负系数。本章新增的是显式的任意固定领口构造、两尺度一致余项和固定基底覆盖上的复制推论；其首阶与高阶费用均按原曲率规范计算。\cite{R09,System,CellSource}

\section{构造结论}\label{ch9:ux6784ux9020ux7ed3ux8bba}

在 \(\operatorname{Ric}=-6g\) 规范下，记 \(P=\int|\operatorname{Ric}^{\circ}|^2\)、 \(J_6=\int|\operatorname{Ric}+6g|^2\)。对每个固定 \(R>0\)，构造普通负单吹起单元 \(g_\varepsilon\)，在 \(t\ge R\) 精确等于复双曲度量，并有 \[
\frac{P(g_\varepsilon)}{8\pi^2}
=1-9\varepsilon^2\log(1/\varepsilon)+O_R(\varepsilon^2),
\] \[
\frac{J_6(g_\varepsilon)}{8\pi^2}
=1-9\varepsilon^2\log(1/\varepsilon)+O_R(\varepsilon^2).
\tag{9.1}\label{eq:r9:1}
\] 负对数项压过全部二次余项，因此充分小的固定正 \(\varepsilon\) 给出严格节省。颈部常模的负系数来自系统原稿；下面保留完整光滑匹配与一致余项推导。

\section{一个明确的光滑度量族}\label{ch9:ux4e00ux4e2aux660eux786eux7684ux5149ux6ed1ux5ea6ux91cfux65cf}

令 \(T=R/2\)。固定 \(\chi\in C^\infty([0,\infty))\)，\(0\le\chi\le1\)，\(\chi=0\) 于 \([0,1]\)，\(\chi=1\) 于 \([2,\infty)\)。例如在 \(1<x<2\) 可取

\[
\chi(x)=\frac{e^{-1/(x-1)}}{e^{-1/(x-1)}+e^{-1/(2-x)}}.
\]

设 \(\epsilon<T^2/16\)。写

\[
S(t)=\sinh t,\quad C(t)=\cosh t,\quad
\mathbf c_*=(9/4,3/4,1/2).
\]

三组轨道半径为

\[
\mathbf a^B=(\sqrt{t^2+\epsilon},\sqrt{t^2+\epsilon},t),
\] \[
\mathbf a^H=(S,SC,S),
\] \[
\mathbf a^N=
\left(\sqrt{S^2+\epsilon}e^{9\epsilon/4},
\sqrt{S^2+\epsilon}\,C e^{3\epsilon/4},
S e^{\epsilon/2}\right).
\]

上标分别表示平直端 Burns 核心、复双曲背景、所选颈部。

令 \(\chi_{\rm in}(t)=\chi(t/\sqrt\epsilon)\)、\(\chi_{\rm out}(t)=1-\chi(t/T)\)。对 \(t>0\) 定义

\[
 a_i(t)=a_i^H(t)\exp\left\{
 \chi_{out}(t)\left[
 \log\frac{a_i^B(t)}{a_i^H(t)}+
 \chi_{in}(t)\log\frac{a_i^N(t)}{a_i^B(t)}
 \right]\right\},
\] \[
 g_\epsilon=dt^2+\sum_{i=1}^3a_i(t)^2\sigma_i^2.
\tag{9.2}\label{eq:r9:2}
\]

所有半径在 \(t>0\) 正。该公式分别给出：

\begin{itemize}
\tightlist
\item
  \(t\le\sqrt\epsilon\)：精确 Burns 核心；
\item
  \(2\sqrt\epsilon\le t\le T\)：精确选定颈部 a\^{}N；
\item
  \(t\ge2T=R\)：精确复双曲背景 a\^{}H。
\end{itemize}

在 \(t=0\) 附近，\(a_1=a_2=\sqrt{t^2+\epsilon}\)、\(a_3=t\)，是普通 Hopf bolt 的标准光滑形式。不存在 bolt 处的截断或分母问题。崩塌 \(\sigma_3\) 圆的全周期给普通 \(\mathcal O(-1)\) 拓扑；外 CH 的 Kähler 纤维为 \(\sigma_2\) 方向，故是与内核 Kähler 方向正交的插入。采用原档案的普通吹起边界标记，手术数保持不变。

各区间之间的 \(\chi\) 在端点平坦，故度量全局 \(C^\infty\)。这不是先构造 \(C^2\) 度量再要求额外光滑化证书。

\section{完整曲率代价的精确校准}\label{ch9:ux5b8cux6574ux66f2ux7387ux4ee3ux4ef7ux7684ux7cbeux786eux6821ux51c6}

令循环指标 (i,j,k)，定义

\[
 u_i=\dot a_i-\frac{a_j^2+a_k^2-a_i^2}{a_ja_k},\qquad
 H_i=u_i+\frac12u_ju_k,\qquad
 \lambda_i=\frac{a_ja_k}{a_i}.
\]

原稿 Cartan 公式给

\[
\frac{P}{8\pi^2}=\int\sum_i\frac{\lambda_i}{4}
 \left(\dot u_i-\frac{2H_i}{\lambda_i}\right)^2dt.
\tag{9.3}\label{eq:r9:3}
\]

定义相对密度

\[
\begin{aligned}
\mathscr L_P={}&\frac{\lambda_1}{4}
(\dot u_1+2H_1/\lambda_1)^2\\
&+\frac{\lambda_2}{4}
(\dot u_2-(2u_2-u_1u_3)/\lambda_2)^2\\
&+\frac{\lambda_3}{4}
(\dot u_3+2H_3/\lambda_3)^2
+\frac{2u_1u_2u_3}{\lambda_2}.
\end{aligned}
\]

逐项展开得到

\[
 \mathscr L_P-\sum_i\frac{\lambda_i}{4}
 (\dot u_i-2H_i/\lambda_i)^2
 =\frac d{dt}(u_1^2+u_3^2+u_1u_2u_3).
\]

bolt 端 \(u=(0,0,-1)\)，外 CH 端 \(u=(0,3\sinh^2t,0)\)。故

\[
\boxed{\frac{P}{8\pi^2}=1+\int_0^R\mathscr L_P\,dt.}
\tag{9.4}\label{eq:r9:4}
\]

这条恒等式保留了整个核心、过渡与外边界，不能以一段颈部积分代替它。

再令

\[
 \Sigma=\sum_i\left(\frac{\dot u_i}{a_i}
 +\frac{2H_i}{a_ja_k}\right),\qquad s=-2\Sigma.
\]

因为 \(dV=2\pi^2a_1a_2a_3\,dt\)，且 \(J_6=P+\frac14\int(s+24)^2\)，故

\[
 \frac{J_6}{8\pi^2}
 =1+\int_0^R\mathscr L_Jdt,\qquad
 \mathscr L_J=\mathscr L_P+
 \frac{a_1a_2a_3}{4}(\Sigma-12)^2.
\tag{9.5}\label{eq:r9:5}
\]

\section{颈部的两尺度解析因子化}\label{ch9:ux9888ux90e8ux7684ux4e24ux5c3aux5ea6ux89e3ux6790ux56e0ux5b50ux5316}

在 \(2\sqrt\epsilon\le t\le T\)，置

\[
 z=\epsilon/t^2,\quad y=t^2,
 \mathfrak s(y)=\frac{\sinh\sqrt y}{\sqrt y},\quad
 \mathfrak c(y)=\cosh\sqrt y.
\]

其中 \(0\le z\le1/4\)、\(0\le y\le T^2\)，且在 \(y=0\) 连续解析延拓。写 \(a_i=tA_i(z,y)\)，则

\[
 A_1=\sqrt{\mathfrak s^2+z}\,e^{zyc_1},\quad
 A_2=\sqrt{\mathfrak s^2+z}\,\mathfrak c\,e^{zyc_2},\quad
 A_3=\mathfrak s\,e^{zyc_3}.
\]

先允许任意固定 \(c=(c_1,c_2,c_3)\)。所有函数在所述实矩形附近实解析且半径为正。以固定 \(\varepsilon\) 微分时，

\[
 t\frac d{dt}=2y\partial_y-2z\partial_z.
\]

因而 \(u_i=U_i(z,y)\)、\(\lambda_i=t\Lambda_i(z,y)\) 均可在该矩形解析表示。

在 \(z=0\) 上是精确 CH，\(U_1=U_3=0\)；在 \(y=0\) 上是缩放后的平直 Burns 数据，

\[
 U_1=U_2=0,\quad U_3=-\frac z{1+z},\quad
 (\Lambda_1,\Lambda_2,\Lambda_3)=(1,1,1+z).
\]

由这两条精确边界校准：

\begin{enumerate}
\def\labelenumi{\arabic{enumi}.}
\tightlist
\item
  \(U_1\) 可被 \(yz\) 整除；\(U_2\) 可被 \(y\) 整除；\(U_3\) 可被 \(z\) 整除。
\item
  定义
\end{enumerate}

\[
\begin{aligned}
B_1&=2(yU_{1,y}-zU_{1,z})+2H_1/\Lambda_1,\\
B_2&=2(yU_{2,y}-zU_{2,z})-(2U_2-U_1U_3)/\Lambda_2,\\
B_3&=2(yU_{3,y}-zU_{3,z})+2H_3/\Lambda_3.
\end{aligned}
\]

每个 \(B_i\) 在 \(z=0\) 与 \(y=0\) 都为零，故 \(B_i=yzb_i\)，\(b_i\) 解析。随后， \(\Sigma\) 可写为 \(y^{-1}B_s\)；\(B_s\) 在 \(y=0\) 为零，而 \(B_s(0,y)=12y\)。因此 \(\Sigma-12=zb_s(z,y)\)，\(b_s\) 解析。

这些因子给出\textbf{完整且无形式余项的表示}

\[
\boxed{\mathscr L_P=\frac{\epsilon^2}{t}F_P(z,y;c),\qquad
\mathscr L_J=\frac{\epsilon^2}{t}F_J(z,y;c),}
\tag{9.6}\label{eq:r9:6}
\]

\(F_P,F_J\) 在整个所述矩形上光滑。由普通均值估计，

\[
 |F_\bullet(z,y;c)-F_\bullet(0,0;c)|\le C_{R,c}(z+y).
\tag{9.7}\label{eq:r9:7}
\]

这一步控制整个颈部，包括下端 \(z=\varepsilon/t^2\) 并不极小的区域；并非将固定 \(t\) 的展开直接积分到 \(t\approx\sqrt\varepsilon\)。

\section{角点系数与严格余项}\label{ch9:ux89d2ux70b9ux7cfbux6570ux4e0eux4e25ux683cux4f59ux9879}

系数可由固定 t 的一阶半径变化计算：

\[
\frac{\partial_\epsilon a_i}{a_i}\Big|_{\epsilon=0}
=\left(\frac1{2\sinh^2t}+c_1,
\frac1{2\sinh^2t}+c_2,c_3\right).
\]

系数计算得到

\[
\begin{aligned}
F_P(0,0;c)={}&11\sum_i c_i^2
 -10(c_1c_2+c_1c_3+c_2c_3)\\
&-37c_1+11c_2+19c_3+59/4,\\
F_J(0,0;c)={}&F_P(0,0;c)+(2c_1+2c_2+2c_3-7)^2/4.
\end{aligned}
\]

令 \(d=c-(9/4,3/4,1/2)\)。它们等价于

\[
F_P=-18+16|d|^2-5(d_1+d_2+d_3)^2,
\] \[
F_J=-18+16|d|^2-4(d_1+d_2+d_3)^2.
\]

两者唯一最小点同为 \(c_*\)，值为 \(-18\)。这个常模最小值原稿已有；此处保留原系数并给出对应二次型。

选择 \(c=c_*\)。由\eqref{eq:r9:6}--\eqref{eq:r9:7}，

\[
\begin{aligned}
\int_{2\sqrt\epsilon}^{T}\mathscr L_\bullet dt
&=-18\epsilon^2\log\frac{T}{2\sqrt\epsilon}
 +O_R\left(\epsilon^2\int_{2\sqrt\epsilon}^{T}
 \frac{t^2+\epsilon/t^2}{t}\,dt\right)\\
&=-9\epsilon^2\log(1/\epsilon)+O_R(\epsilon^2).
\end{aligned}
\]

误差积分确实有与 \(\varepsilon\) 无关的上界：

\[
\int_{2\sqrt\epsilon}^{T}t\,dt\le T^2/2,\qquad
\int_{2\sqrt\epsilon}^{T}\epsilon t^{-3}dt\le1/8.
\]

因此余项为 \(O_R(\varepsilon^2)\)，严格低于决定符号的 \(\varepsilon^2\log(1/\varepsilon)\) 阶。

\section{\texorpdfstring{内核与外过渡都只支付 \(O_R(\epsilon^2)\)}{内核与外过渡都只支付 O\_R(\textbackslash epsilon\^{}2)}}\label{ch9:ux5185ux6838ux4e0eux5916ux8fc7ux6e21ux90fdux53eaux652fux4ed8-o_repsilon2}

\textbf{内核。} \(t\le\sqrt\epsilon\) 上精确 Burns，\(\mathscr L_P=0\)。标量项的归一化积分恰为 \(27\epsilon^2\)。

在 \(\sqrt\epsilon\le t\le2\sqrt\epsilon\) 上令 \(t=\sqrt\epsilon\,r\)。缩放后度量 \(\epsilon^{-1}g_\epsilon\) 是固定 Burns 核心在 \(1\le r\le2\) 上的 \(C^2\) 扰动，大小 \(O(\epsilon)\)。其 \(u_1,u_2=O(\epsilon)\)，而 Burns 的三个校准平方都为零。于是校准平方与三次项的积分全部为 \(O(\epsilon^2)\)。标量曲率在原尺度为 \(O(1)\)，该区域体积为 \(O(\epsilon^2)\)，故 \(J_6-P\) 也为 \(O(\epsilon^2)\)。内区临界处没有裸的一阶 \(\varepsilon\) 费用。

\textbf{外过渡。} \(T\le t\le2T\) 是固定紧环带。\(g_\epsilon-g_H=O_R(\epsilon)\) 于 \(C^2\)，\(g_H\) 为 \(\operatorname{Ric}=-6g\)。故 \(|\operatorname{Ric}+6g|=O_R(\epsilon)\)，其积分费用 \(O_R(\epsilon^2)\)。同样 P=\(O_R(\epsilon^2)\)。也可以直接在\eqref{eq:r9:4}中看到：\(u_1,u_3=O_R(\epsilon)\)，每个校准平方为 \(O_R(\epsilon^2)\)，三次项 \(u_1u_2u_3\) 也为 \(O_R(\epsilon^2)\)。

合并三段证明\eqref{eq:r9:1}。光滑性、正定性、普通吹起拓扑与完整标记领口均已在构造中满足。

\section{归一化与每个预定覆盖塔}\label{ch9:ux5f52ux4e00ux5316ux4e0eux6bcfux4e2aux9884ux5b9aux8986ux76d6ux5854}

缩放 \(\widetilde g=6g\) 后，背景为 \(\operatorname{Ric}=-\widetilde g\)，且 \(J_1(\widetilde g)=J_6(g)\)。若用异常面积参数 \(x=\operatorname{Area}(E)/\pi\)，则 \(x=6\varepsilon\)，故

\[
\frac{J_1}{8\pi^2}
=1-\frac{x^2}{4}\log(1/x)+O_R(x^2).
\]

考虑固定紧复双曲曲面 \(X_0\) 及任意有限非分歧覆盖序列 \(X_j\to X_0\)，度数 \(d_j\to\infty\)。在 \(X_0\) 上选择小嵌入球并留足缓冲；\eqref{eq:r9:1}提供适合这个球的固定光滑单元 C，满足

\[
 J_1(C)=8\pi^2(1-\delta_{X_0}),\qquad\delta_{X_0}>0.
\]

每个覆盖含 d\_j 个互不相交的提升球。对原比例

\[
 Q_j=k_j+4\ell_j,\quad h_j=Q_j/A_j\to0,\quad
 k_j/Q_j\to2/9,\quad \ell_j/Q_j\to7/36,
\]

因 A\_j=d\_jA\_0，k\_j+2\(\ell\)\_j=o(d\_j)，所以充分大指标下有足够的提升球放置所有固定吹起单元和缩小的手柄部分。使用原供给及复制公式，得到

\[
\boxed{
\limsup_j\frac{\alpha(M_j)}{h_j}
\le\frac{28}{9}+\frac{J_1(C)}{9\pi^2}
=4-\frac89\delta_{X_0}<4.
}
\tag{9.8}\label{eq:r9:8}
\]

此处不要求先将给定覆盖塔替换为容纳大半径数值单元的特殊覆盖塔。常数 \(\delta\)\_\{X0\} 可依赖固定基底和选定球；没有证明它对所有不同基底统一。

\textbf{极限次序至关重要。} 先固定球，再选取一个足够小但固定的 \(\varepsilon\)，然后令 j\(\to\)\(\infty\)。若让 \(\varepsilon\)\_j\(\to\)0，节省也趋零，便不能据此获得严格小于4的固定首项界。

\begin{thebibliography}{P99}
\small
\bibitem[R09]{R09} 本会话研究原稿，\emph{任意复双曲领口中的解析低成本构造}，2026-10-07。\path{O11_Arbitrary_Collar_Logarithmic_Construction.md}。

\bibitem[P08]{System} S1003A01、S1003A02，\emph{Mathematical Synthesis}，2026-10-07；\path{20261007_Einstein_Gaps_Two_Sessions(1).zip} 中 \path{01_Mathematical_Synthesis.md}。

\bibitem[P09]{CellSource} S1003A01 Step10；S1003A02 Step11；\emph{Theorem package v1}，完整三轴单元、颈部二次型与严格次四恢复。

\end{thebibliography}
\end{document}
